\section{The essay-scorer analogy}
\label{sec:essay-evaluator}

Educational and psychometric measurement draws a line that grammaticality research often blurs: the scored aren't the scorers. Essay scorers, oral-proficiency raters, and constructed-response markers can introduce error, bias, drift, and disagreement, and assessment researchers take that seriously. None of it converts the scorers into participants alongside the students. The same holds inside the research process: a colleague or reviewer who judges whether an author's example is acceptable, or whether a claimed contrast holds, is assessing the materials, not thereby enrolling as a subject.

What makes the analogy bite is that this scoring isn't casual. Constructed-response programmes train raters against rubrics, double-score responses, monitor severity and drift, adjudicate disagreements, and model rater effects statistically \citep{mccaffrey2022scoring,aera2014standards}. The scoring is about as systematic, multi-rater, and quantitatively scrutinized as human judgment gets. The raters are still part of the measurement apparatus, not its object.

That ordering matters, because the usual worry about grammaticality judgments runs the other way. The intuition is that once judgment collection becomes systematic, with many judges, fixed materials, and reliability statistics, it must be participant research. The essay-scorer case shows the inference backwards. Systematicity, multiple raters, and reliability analysis are signatures of a measurement procedure, not marks of a study about the judges.

The sharpest form of the worry presses on reliability itself. If a study reports how far several judges agree, isn't it then reporting facts about those people? It's reporting a property of the measurement, not a finding about the judges. Inter-rater agreement is what licenses the step from the judges' verdicts to the status of the object, much as a rater facet functions in psychometric models \citep{aera2014standards}. The statistic characterizes the instrument; the object stays what the study is about.

The line can move, though. An aggregate agreement coefficient characterizes the procedure, but a study that models rater severity by training, dialect, gender, or theoretical orientation has stopped treating rater variation as a nuisance facet and made the raters part of what it studies.

The analogy buys neither infallibility nor exemption from scrutiny. It relocates the scrutiny. If raters can be unreliable, the response is to train, blind, rotate, calibrate, adjudicate, report: the standards that govern scoring, not those that govern participant sampling \citep{aera2014standards}. An expert grammaticality judge deciding whether an item behaves as a theory implies is doing the rater's job, and should be held to the rater's standards.

The analogy isn't exact, and the gap cuts in a useful direction. Essay scorers normally judge performances produced independently of them, whereas grammaticality evaluators often judge examples constructed for the argument they support. That raises the need for disclosure, independence, and, where feasible, blinding to the prediction. It changes the controls the role requires, not the role itself.

A deeper gap runs the other way. An essay's merits are fixed independently of the scorer, but a sentence's grammaticality supervenes on the community's competence, the faculty the judgment draws on. The object is a constituted, community-level kind, accessed through the evaluator's competence rather than grounded in it. The from/about line still holds, since the report isn't about this judge, but the analogy's object-independence is partly borrowed.
